% ===========================================================================
% Séquence — candidature, analyse de CV, notation, classement.
%
% Tracé sur le code réel (2026-09-22) :
%   api/v1/endpoints/recruitment.py::apply_to_job  (ligne 237, tâche ligne 304)
%   agents/scoring_agent.py::run_analysis_task / analyze_application
%   services/{cv_parser,nlp_service,semantic_service,scoring_service}.py
%
% Le point de conception à faire ressortir : la réponse 201 part AVANT que
% l'analyse ne commence. Le candidat n'attend pas l'extraction NLP, qui dure
% plusieurs secondes. La tâche de fond ouvre sa PROPRE session de base, car
% celle de la requête est fermée dès la réponse envoyée.
% ===========================================================================
\begin{tikzpicture}[x=1mm, y=-1mm]

  \def\lifelen{150}

  % ─── Participants ───────────────────────────────────────────────────────
  \node[pqhead] (c)   at (0,0)   {Candidat};
  \node[pqhead] (api) at (34,0)  {API\\\texttt{/apply}};
  \node[pqhead] (ag)  at (72,0)  {Agent de notation\\{\scriptsize tâche de fond}};
  \node[pqhead] (sv)  at (105,0) {Services\\{\scriptsize CV, NLP, score}};
  \node[pqhead] (db)  at (136,0) {PostgreSQL};

  \foreach \n in {c,api,ag,sv,db}
    \draw[pqlifeline] (\n.south) -- ++(0,\lifelen);

  % ═══ Phase 1 — dépôt synchrone ══════════════════════════════════════════
  \fill[pqactivation] (32.5,16) rectangle (35.5,52);

  \draw[pqmsg] (0,16) -- node[pqmsglbl] {\texttt{POST} candidature + CV} (32.5,16);

  % Auto-appel : boucle courte sur la ligne de vie, plutôt qu'une flèche de
  % longueur nulle — illisible au rendu.
  \draw[pqmsg] (35.5,24) -- ++(9,0) -- ++(0,7) -- (35.5,31);
  \node[pqnote, anchor=west] at (46,27.5)
    {valide le fichier — liste blanche de types, plafond 5\,Mo};

  \draw[pqmsg]    (35.5,39) -- node[pqmsglbl] {enregistre la candidature} (136,39);
  \draw[pqreturn] (136,45)  -- (35.5,45);

  \draw[pqmsg] (35.5,58) -- node[pqmsglbl] {\textbf{planifie} l'analyse} (72,58);
  \node[pqnote, anchor=north west] at (74,61) {tâche de fond, session dédiée};

  \draw[pqreturn] (32.5,70) -- node[pqmsglbl] {\texttt{201 Created}} (0,70);

  \node[pqwarn, anchor=north west, text width=60mm] at (-4,76)
    {\textbf{La réponse part avant l'analyse.} Le candidat n'attend pas
     l'extraction, qui prend plusieurs secondes.};

  % ═══ Séparateur de phase ════════════════════════════════════════════════
  \draw[draw=piqbitExt, dash pattern=on 2pt off 2pt] (-10,100) -- (148,100);
  \node[anchor=west, font=\scriptsize\itshape, text=piqbitExt,
        fill=white, inner sep=2pt] at (-10,100)
    {traitement asynchrone — la session de la requête est déjà fermée};

  % ═══ Phase 2 — analyse asynchrone ═══════════════════════════════════════
  \fill[pqactivation] (70.5,108) rectangle (73.5,146);

  \draw[pqmsg]    (73.5,112) -- node[pqmsglbl] {extraire le texte du CV} (105,112);
  \draw[pqmsg]    (73.5,120) -- node[pqmsglbl] {extraire compétences, expérience} (105,120);
  \draw[pqreturn] (105,126)  -- (73.5,126);

  \draw[pqmsg] (73.5,133) -- node[pqmsglbl] {indexer les vecteurs} (105,133);
  \node[pqnote, anchor=west] at (107,133) {\emph{best-effort}};

  \draw[pqmsg] (73.5,141) -- node[pqmsglbl] {calculer le score, classer} (105,141);
  \draw[pqmsg] (105,148)  -- node[pqmsglbl] {\texttt{COMMIT}} (136,148);

  \node[pqnote, anchor=north west, text width=76mm] at (-4,154)
    {Un échec d'indexation sémantique n'annule pas la notation : l'analyse
     reste exploitable sans recherche vectorielle.};

\end{tikzpicture}
